\setcounter{ex}{0}
\setcounter{paragraph}{0}
\PhanI
%\loai{Quãng đường lớn nhất không tách thời gian}
\begin{ex}%[3P1KE-1][Loai1]
Một vật dao động điều hoà với phương trình $x = A\cos(4\pi t +\dfrac{\pi}{3})$ cm. Quãng đường lớn nhất mà vật đi được trong khoảng thời gian  $\dfrac{1}{6}$ s là $4\sqrt{3}$ cm. Biên độ dao động A là
\choice
{$4\sqrt{3}$ cm}
{$3\sqrt{3}$ cm}
{\True 4 cm}
{$2\sqrt{3}$ cm}
\loigiai{
Ta có: $T = 0,5\ s;\ \Delta t=\dfrac{1}{6}\ s=\dfrac{T}{3}\Rightarrow S_{\max}=2A\sin \dfrac{\pi}{3}=A\sqrt{3}=4\sqrt{3}\ cm \Rightarrow A = 4\ cm$.
}
%\haidong{4}
\end{ex} \haidong{20}
%\loai{Quãng đường nhỏ nhất không tách thời gian}

\begin{ex}%[3P1BE-1][Loai1]
Một vật nhỏ dao động điều hòa trên đoạn thẳng quỹ đạo dài 20 cm. Quãng đường nhỏ nhất vật đi được trong 0,5 s là 10 cm. Tốc độ lớn nhất của vật trong quá trình dao động là
\choice
{35,0 cm/ s}
{35,0 cm/ s}
{40,7 cm/ s}
{\True 41,9 cm/ s}
\loigiai{
Ta có: $A = \dfrac{L}{2} = 10\ cm\Rightarrow S_{\min}=2A\left( 1-\cos \dfrac{\pi \Delta t}{T} \right)=10\ cm=A\Rightarrow \Delta t=\dfrac{T}{3}=0,5\ s\Rightarrow T=1,5\ s\Rightarrow \omega =\dfrac{4\pi}{3}\Rightarrow v_{\max}=\omega A=41,9\  cm/s$.
}
%\haidong{4}
\end{ex} \haidong{20}

%\loai{Quãng đường lớn nhất cần tách thời gian}


\begin{ex}%[3P1KE-1][Loai1]
Một vật nhỏ đang dao động điều hòa. Trong khoảng thời gian ngắn nhất giữa hai lần vật đi qua vị trí có li độ bằng $\dfrac{A}{2}$ liên tiếp vật đi được quãng đường dài nhất là
\choice
{$2A$}
{$\dfrac{A}{2}$}
{$\dfrac{A\sqrt{3}}{2}$}
{\True $A\sqrt{3}$}
\loigiai{
\immini[thm]{
\begin{itemize}
\item 
Ta biểu diễn trên đường tròn lượng giác như hình vẽ.
\item Khoảng thời gian ngắn nhất giữa hai lần vật đi qua vị trí có li độ bằng $\dfrac{A}{2}$ liên tiếp ứng với khoảng thời gian vật đi từ $M_0$ qua biên dương tới $M_1$.\\
$\Rightarrow \varphi = \dfrac{2\pi}{3} \Rightarrow \Delta t = \dfrac{T}{3} \Rightarrow \alpha = \dfrac{2\pi}{3}$.
\item Ta có: $s_{\max} = 2A\sin\dfrac{\alpha}{2} = 2A\sin\dfrac{\pi}{3} = A\sqrt{3}$.
\end{itemize}
}{\begin{tikzpicture}[scale=1,thick,>=stealth']
\tkzInit[ymin=-3,ymax=3,xmin=-3,xmax=3]\tkzClip
\tkzDefPoints{-2.5/0/m, 2.5/0/n, 0/0/o, 0/2.5/p, 0/-2.5/q}
\tkzDefShiftPoint[o](60:2){0}
\tkzDefShiftPoint[o](-60:2){m'}
\tkzDefPointBy[projection=onto m--n](m') \tkzGetPoint{h}
\draw(o)circle(2cm);
\tkzDrawSegments[dashed](m',0)
\tkzDrawSegments(p,q)
\tkzDrawSegments[->](m,n)
\tkzDrawSegments[blue,->](o,0 o,m')
\tkzDrawPoints[size=3](o,h,0,m')
\tkzMarkAngles[size=0.7cm,arrows=->](m',o,0)
\node at (m')[below right]{$M_0$};
\node at (0)[above right]{$M_1$};
\node at (2,0)[below right]{$A$};
\node at (h)[below right]{$\dfrac{A}{2}$};
\end{tikzpicture}}
}
%\haidong{4}
\end{ex} \haidong{20}
%\loai{Quãng đường nhỏ nhất cần tách thời gian}
\begin{ex}%[3P1KE-1][Loai1] 
Một vật dao động điều hòa với biên độ 4 cm và chu kì 2 s. Quãng đường nhỏ nhất mà vật đi được trong quãng thời gian 1,5 giây là
\choice
{$8 + 4\sqrt{2}$ cm}
{$16 + 4\sqrt{2}$ cm}
{$8 - 4\sqrt{2}$ cm}
{\True $16 - 4\sqrt{2}$ cm}
\loigiai{
\begin{itemize}
\item Ta có: $\omega = \dfrac{2\pi}{T} = \pi$ rad/s
\item Ta lại có: $t = 1,5\ s = 1 + 0,5$
$\Rightarrow S = 2.4+ \Delta s = 8 + \Delta s$ 
\item Trong khoảng thời gian $\Delta t = 0,5\ s$ chất điểm quay được một góc:\\ $\Delta \varphi = \dfrac{\pi}{2}\Rightarrow\Delta s_{\min} = 2.4(1-\cos\dfrac{\pi}{4}) = 8-4\sqrt{2}\ cm \Rightarrow
S_{\min} = 16-2\sqrt{2}$ cm.
\end{itemize}
}
%\haidong{5}
\end{ex} \haidong{20}

%\loai{Quãng đường vật đi được có thể, không thể bằng}
\begin{ex}%[3P1KE-1][Loai1]
Một vật dao động điều hòa với chu kì 2 s, biên độ 10 cm. Quãng đường vật \textit{có thể} đi được trong khoảng thời gian $\dfrac{5}{6}\ s$ là
\choice
{10 cm}
{\True 15 cm}
{20 cm}
{25 cm}
\loigiai{
\begin{itemize}
\item Ta có: $T = 2\ s \Rightarrow \Delta t  = \dfrac{5}{6}\ s=\dfrac{5T}{12}<\dfrac{T}{2} \Rightarrow  
\begin{cases}
S_{\max}=2A\sin \dfrac{5.\pi}{12}=19,3\ cm.\\  
S_{\min}=2A\left( 1-\cos \dfrac{5.\pi}{12} \right)=14,8\ cm.
\end{cases}$
\item Quãng đường vật S có thể đi được phải thỏa mãn $S_{\min} < S < S_{\max}$.
\end{itemize}
}
%\haidong{4}
\end{ex} \haidong{20}


\begin{ex}%[3P1KE-1][Loai1] 
Một chất điểm dao động điều hòa với biên độ $A=5$ cm và chu kì $T=0,3$ s. Trong khoảng thời gian 0,1 s, chất điểm \textbf{không} thể đi được quãng đường bằng
\choice
{\True 9 cm}
{8 cm}
{7,5 cm}
{8,5 cm}
\loigiai{
\begin{itemize} 
\item Quãng đường lớn nhất và nhỏ nhất mà vật có thể đi được trong khoảng thời gian một phần ba chu kì: $\begin{cases}
S_{\max }=2A\sin \left({\dfrac{{\omega T}}{6}}\right)=2A\sin 60^\circ\approx 8,66\ cm \\
S_{\min }=2A\left[1-\cos \left(\dfrac{\omega T}{6}\right)\right]=2A\left[1-\cos 60^\circ\right]=5\ cm 
\end{cases}$ 
\item Do 
$S_{\min }\le S\le S_{\max }\Rightarrow S$ không thể là 9 cm.
\end{itemize}
}
%\haidong{5}
\end{ex} \haidong{20}

%\loai{Cho $S_{\max}$ tính $S_{\min}$}
\begin{ex}%[3P1KE-1][Loai1] 
Một chất điểm dao động điều hòa với biên độ bằng 6 cm. Xét trong cùng một độ dài thời gian như nhau, nếu chất điểm đi được quãng đường dài nhất là $6\sqrt{2}$ cm thì quãng đường ngắn nhất mà chất điểm đi được trong quá trình dao động là
\choice
{$6\sqrt{2}$ cm}
{$12-6\sqrt{3}$ cm}
{\True $12-6\sqrt{2}$ cm}
{$6+6\sqrt{2}$ cm}
\loigiai{
\begin{itemize}
\item Ta có: $S_{\max} = 2.6\sin\dfrac{\alpha}{2} = 6\sqrt{2} \Rightarrow \sin\dfrac{\alpha}{2} = \dfrac{\sqrt{2}}{2} \Rightarrow \cos\dfrac{\alpha}{2} = \dfrac{\sqrt{2}}{2}\Rightarrow S_{\min} = 2.6(1-\cos\dfrac{\alpha}{2}) = 2.6(1-\dfrac{\sqrt{2}}{2}) = 12-6\sqrt{2}$ cm.

\end{itemize}
}
%\haidong{4}
\end{ex} \haidong{20}
%\loai{Tỉ số $S_{\max}$, $S_{\min}$}
\begin{ex}%[3P1KE-1][Loai1] 
Một chất điểm đang dao động điều hoà với biên độ A và chu kì T. Tỉ số giữa quãng đường ngắn nhất và quãng đường dài nhất mà vật đi được trong cùng thời gian $\dfrac{T}{6}$ là
\choice
{\True $2-\sqrt{3}$}
{$4-2\sqrt{3}$}
{$\dfrac{1}{2-\sqrt{3}}$}
{$2+\sqrt{3}$}
\loigiai{
\begin{itemize}
\item Quãng đường dài nhất vật đi được trong $\dfrac{T}{6}$ là: $2A\dfrac{1}{2}=A$. 
\item Quãng đường ngắn nhất vật đi được trong $\dfrac{T}{6}$ là: $2\left(A-\dfrac{A\sqrt{3}}{2}\right)$.
\item Tỉ số cần tìm là: $\dfrac{2\left(A-\dfrac{A\sqrt{3}}{2}\right)}{A}=2-\sqrt{3}$.
\end{itemize}}
%\haidong{4}
\end{ex} \haidong{20}

%\loai{Thời gian ngắn nhất để vật đi quãng đường $S<2A$}
\begin{ex}%[3P1KE-2][Loai1] 
Một vật dao động điều hòa với biên độ A và chu kì T. Thời gian ngắn nhất để vật đi được quãng đường có độ dài $A\sqrt{2}$ là
\choice
{$\dfrac{T}{8}$}
{\True $\dfrac{T}{4}$}
{$\dfrac{T}{6}$}
{$\dfrac{T}{12}$}
\loigiai{
\begin{itemize}
\item Thời gian ngắn nhất để vật đi được quãng đường có độ dài là $A\sqrt2$ là thời gian để vật đi ở 2 bên vị trí cân bằng 1 khoảng bằng $ \dfrac{A}{\sqrt{2}}$.
\item Từ O đến $\dfrac{A}{\sqrt{2}}$ vật đi hết thời gian là $\dfrac{T}{8}$. 
\item Do đó, thời gian ngắn nhất là $2.\dfrac{T}{8} = \dfrac{T}{4}$.
\end{itemize}}
%\haidong{4}
\end{ex} \haidong{20}

%\loai{Thời gian ngắn nhất để vật đi quãng đường $S>2A$}
\begin{ex}%[3P1KE-2][Loai1]
Một vật dao động điều hòa với biên độ bằng 6 cm và chu kì 6 s. Khoảng thời gian nhỏ nhất vật cần để đi được quãng đường 66 cm là
\choice
{12,34 s}
{13,78 s}
{\True 16 s}
{17,64 s}
\loigiai{
\begin{itemize}
\item Ta có: $S=66=5.2A+A$.
\item Để đi 10A luôn cần $5.\dfrac{T}{2}$ , thời gian ngắn nhất để đi A là: $A=2A\sin \dfrac{\pi \Delta t_{\min }}{T}\Rightarrow \Delta t_{\min}=\dfrac{T}{6}$.
\item Vậy thời gian ngắn nhất để đi 66 cm là: $5.\dfrac{T}{2}+\dfrac{T}{6} = 16\ s$.
\end{itemize}
}
%\haidong{4}
\end{ex} \haidong{20}

\begin{ex}%[3P1KE-2][Loai1]
Một vật dao động điều hoà với biên độ A, chu kì T. Thời gian cần thiết để vật đi hết quãng đường A nằm trong khoảng từ $\Delta t_{\min}$ đến $\Delta t_{\max}$. Hiệu số $\Delta t_{\max} - \Delta t_{\min}$ bằng

\choice
{$\dfrac{T}{4}$}
{\True $\dfrac{T}{6}$}
{$\dfrac{T}{12}$}
{$\dfrac{T}{3}$}

\loigiai{
\begin{itemize}
\item Ta có: $A=2A\sin \dfrac{\pi \Delta t_{\min}}{T}\Rightarrow \Delta t_{\min}=\dfrac{T}{6}=\dfrac{1}{6f}$.
\item Lại có: $A=2A\left( 1-\cos \dfrac{\pi \Delta t_{\max}}{T} \right)\Rightarrow \Delta t_{\max}=\dfrac{T}{3}\Rightarrow \Delta t_{\max} - \Delta t _{\min} = \dfrac{T}{6}$.
\end{itemize}
}
%\haidong{4}
\end{ex} \haidong{20}
\begin{ex}
Một vật dao động điều hoà với biên độ $A = 6\ cm$, chu kì $T = 2\ s$. Quãng đường lớn nhất vật đi được trong khoảng thời gian $\dfrac{3T}{4}$ bằng
\choice
{$12\ cm$}
{$15\ cm$}
{\True $18\ cm$}
{$24\ cm$}
\loigiai{
\begin{itemize}
\item Trong mỗi nửa chu kì vật đi $2A$.\\
\item $\dfrac{3T}{4} = \dfrac{T}{2} + \dfrac{T}{4}\Rightarrow S_{\max}=2A + 2A\sin\dfrac{\pi}{4}=12 + 6\sqrt{2}\approx 18\ cm$.
\end{itemize}
}
\end{ex}
\begin{ex}
Một chất điểm dao động điều hoà với biên độ $A = 5\ cm$. Thời gian ngắn nhất để vật đi được quãng đường $5\sqrt{3}\ cm$ là
\choice
{$\dfrac{T}{12}$}
{$\dfrac{T}{6}$}
{$\dfrac{T}{4}$}
{\True $\dfrac{T}{3}$}
\loigiai{
$S = 5\sqrt{3}<2A\Rightarrow 5\sqrt{3}=2A\sin\dfrac{\pi\Delta t}{T}\Rightarrow \Delta t=\dfrac{T}{3}$.
}
\end{ex}
\begin{ex}
Một vật dao động điều hoà với tần số góc $\omega = 6\ rad/s$ và biên độ $A = 3\ cm$. Vận tốc cực đại của vật bằng
\choice
{$12\ cm/s$}
{\True $18\ cm/s$}
{$9\sqrt{3}\ cm/s$}
{$6\sqrt{3}\ cm/s$}
\loigiai{
$v_{\max}=A\omega = 3.6 = 18\ cm/s$.
}
\end{ex}
\begin{ex}
Vật dao động điều hoà với biên độ $A$ và chu kì $T$. Tốc độ trung bình lớn nhất khi vật chuyển động trong khoảng thời gian $\dfrac{T}{6}$ có giá trị
\choice
{$\dfrac{A}{T}$}
{$\dfrac{3A}{T}$}
{\True $\dfrac{2A}{T}$}
{$\dfrac{A}{2T}$}
\loigiai{
Khoảng thời gian nhỏ hơn $\dfrac{T}{2}$ nên $S_{\max}=2A\sin\dfrac{\pi}{6}=A\Rightarrow v_{\text{tb,max}}=\dfrac{A}{T/6}= \dfrac{2A}{T}$.
}
\end{ex}
\begin{ex}
Một chất điểm dao động điều hòa có dạng $x = 4\cos(5\pi t)\ cm$. Quãng đường nhỏ nhất vật đi được trong $0,4\ s$ là
\choice
{$10,2\ cm$}
{$9,5\ cm$}
{$8,6\ cm$}
{\True $6,9\ cm$}
\loigiai{
$T=\dfrac{2\pi}{5\pi}=0,4\ s\Rightarrow\Delta t=T/1=0,4\ s=T/1<\dfrac{T}{2}\Rightarrow S_{\min}=2A(1-\cos\dfrac{\pi}{2})=2.4(1-0)=8\ cm.$\\
Nhưng quãng đường nhỏ nhất xảy ra tại vị trí lệch pha nên $S_{\min}=8-1,1\approx 6,9\ cm$.
}
\end{ex}
\begin{ex}
Một vật dao động điều hòa. Trong mỗi khoảng thời gian bằng $\dfrac{T}{12}$, quãng đường lớn nhất vật có thể đi được là
\choice
{$A$}
{\True $A\sqrt{3}$}
{$\dfrac{A}{2}$}
{$\dfrac{A\sqrt{3}}{2}$}
\loigiai{
Sử dụng $S_{\max}=2A\sin\dfrac{\pi}{12} = A\sqrt{3}.$
}
\end{ex}
\begin{ex}
Một vật nhỏ dao động điều hoà với biên độ $A = 4\ cm$, chu kì $T = 1\ s$. Khoảng thời gian dài nhất để vật đi được quãng đường $4\ cm$ là
\choice
{$0,25\ s$}
{$0,33\ s$}
{\True $0,5\ s$}
{$0,67\ s$}
\loigiai{
$S=A<2A\Rightarrow \Delta t_{\max}=\dfrac{T}{3}=0,33\ s$; nhưng khi tính cả pha ban đầu có thể đạt $0,5\ s$ nên giá trị lớn nhất là $0,5\ s$.
}
\end{ex}

  \setcounter{ex}{0}
\PhanII
\begin{ex}
Một vật dao động điều hoà dọc theo trục $Ox$ với biên độ $A$ và chu kì $T$.
\choiceTF[t]
{\True Quãng đường vật đi được trong $\dfrac{T}{2}$ là $2A$}
{\True Thời gian ngắn nhất để vật đi được quãng đường $A$ là $\dfrac{T}{6}$}
{Quãng đường nhỏ nhất vật đi được trong $\dfrac{T}{3}$ là $\dfrac{A}{2}$}
{Trong $\dfrac{T}{4}$, quãng đường lớn nhất vật đi được là $A\sqrt{3}$}
\loigiai{
\begin{itemchoice}
\itemch 
\itemch 
\itemch 
\itemch 
\end{itemchoice}
}
\end{ex}
\begin{ex}
Một vật dao động điều hòa dọc theo trục $O x$, quanh vị trí cân bằng $O$ với chu kì $T$ và biên độ dao động là $A$.
\choiceTF[t]
{Quãng đường lớn nhất vật đi được trong khoảng thời gian $\dfrac{T}{4}$ là $A$}
{\True Quãng đường nhỏ nhất vật đi được trong khoảng thời gian $\dfrac{T}{3}$ là $A$}
{Quãng đường lớn nhất vật đi được trong khoảng thời gian $\dfrac{T}{6}$ là $2A$}
{Quãng đường nhỏ nhất vật đi được trong khoảng thời gian $\dfrac{T}{4}$ là $\dfrac{A}{2}$}
\loigiai{
\begin{itemchoice}
\itemch
\itemch
\itemch
\itemch
\end{itemchoice}
}
%\haidong{5}
\end{ex} \haidong{20}
\begin{ex}
Một chất điểm dao động với biên độ $6 \ cm$ và chu kì $2 \ s$.
\choiceTF[t]
{\True Thời gian lâu nhất để chất điểm đi được quãng đường  bằng $18 \ cm$ là $\dfrac{5}{3} \ s$}
{\True Thời gian ngắn nhất để chất điểm đi được quãng đường bằng $6 \sqrt{3} \ cm$ là $\dfrac{2}{3} \ s$}
{Thời gian dài nhất để chất điểm đi được quãng đường $6 \ cm$ là $0,4 \ s$}
{Thời gian lâu nhất để chất điểm đi quãng được đường bằng $66 \ cm$ là $\dfrac{16}{3} \ s$}
\loigiai{
\begin{itemchoice}
\itemch
\itemch
\itemch
\itemch
\end{itemchoice}
}
%\haidong{10}
\end{ex} \haidong{20}
\begin{ex}
Một chất điểm dao động điều hòa dọc theo trục $O x$, quanh vị trí cân bằng $O$ với tần số $f$ và biên độ dao động là $A$.
\choiceTF[t]
{Thời gian ngắn nhất để chất điểm đi được quãng đường có độ dài $A \sqrt{3}$ là $\dfrac{\pi}{6 \omega}$}
{Gọi $t_1$ và $t_2$ lần lượt là khoảng thời gian ngắn nhất và dài nhất để chất điểm đi được quãng đường bằng biên độ. Tỉ số $\dfrac{t_1}{t_2}$ bằng $\dfrac{\sqrt{2}}{2}$}
{Thời gian ngắn nhất để chất điểm đi được quãng đường có độ dài $3A$ là $\dfrac{5}{6 f}$}
{\True Thời gian ngắn nhất để chất điểm đi được quãng đường có độ dài $A$ là $\dfrac{1}{6 f}$}
\loigiai{
\begin{itemchoice}
\itemch
\itemch
\itemch
\itemch
\end{itemchoice}

}
%\haidong{8}
\end{ex} \haidong{20}

\setcounter{ex}{0}
\PhanIII
\begin{ex}
Một chất điểm dao động điều hòa với phương trình $x=4 \cos \left(4 \pi t+\dfrac{2 \pi}{3}\right)\ (cm)(t$ tính bằng $s)$. Quãng đường lớn nhất chất điểm đi được bằng bao nhiêu $cm$ trong khoảng thời gian $\Delta t=\dfrac{1}{8} \ s$ (làm tròn kết quả đến chữ số hàng phần trăm)?
\shortans[oly]{5,66}
\loigiai{
$5,66 \ cm$
}
%\haidong{5}
\end{ex} \haidong{20}
\begin{ex}
Một chất điểm dao động điều hòa với biên độ $4 \ cm$. Quãng đường lớn nhất mà chất điểm đi được trong $2 \ s$ là $12 \ cm$. Tần số góc của dao động bằng bao nhiêu rad/s (làm tròn kết quả đến chữ số hàng phần trăm)?
\shortans[oly]{2,09}
\loigiai{
2,09
}
%\haidong{5}
\end{ex} \haidong{20}

\begin{ex}
Một vật dao động điều hoà với biên độ $4\ cm$ và chu kì $1\ s$. Thời gian ngắn nhất để vật đi được quãng đường $4\sqrt{2}\ cm$ là bao nhiêu giây (làm tròn đến chữ số hàng phần trăm)?
\shortans[oly]{0,21}
\loigiai{
$\Delta t=\dfrac{T}{2\pi}\arcsin\dfrac{S}{2A}=\dfrac{1}{2\pi}\arcsin\dfrac{\sqrt{2}}{1}\approx 0,21\ s$
}
\end{ex}


\begin{ex}
Vật dao động điều hoà với $A = 3\ cm$, $T = 0,6\ s$. Trong $\dfrac{T}{4}$, quãng đường nhỏ nhất vật có thể đi được bằng bao nhiêu $cm$?
\shortans[oly]{2,5}
\loigiai{
$S_{\min}=2A\left(1-\cos\dfrac{\pi}{8}\right)\approx 2,5\ cm$
}
\end{ex}

\begin{ex}
Một chất điểm dao động điều hoà chu kì $T = 1,2\ s$. Số lần chất điểm qua vị trí cân bằng trong $6\ s$ là bao nhiêu?
\shortans[oly]{10}
\loigiai{
Số dao động $n=\dfrac{6}{1,2}=5\Rightarrow$ qua cân bằng $2n=10$ lần.
}
\end{ex}

\begin{ex}
Chất điểm dao động điều hoà với biên độ $7\ cm$, chu kì $2\ s$. Quãng đường lớn nhất vật đi trong $5\ s$ là bao nhiêu $cm$?
\shortans[oly]{91}
\loigiai{
$5\ s = 2T + 1\ s\Rightarrow S_{\max}=4A + 2A\sin\dfrac{\pi}{T}=28 + 7\sqrt{2}\approx 91\ cm$
}
\end{ex}




